
@@M0@@ qui
envoie @@M1@@ sur @@M2@@, et @@M3@@ sur @@M4@@~; d'après la remarque précédente, l'élément obtenu ne dépend pas de l'entier @@M5@@ tel que @@M6@@ pour @@M7@@. En utilisant le même homomorphisme, on voit que la relation (5.3.1) est vérifiée par @@M8@@ et @@M9@@.

5.4. Considérons le @@M10@@-anneau engendré sur @@M11@@ par les indéterminées @@M12@@, soumises aux relations @@M13@@ pour @@M14@@. De (5.3.1) résulte :
\begin{align*}
  \lambda^p(N,x+y)\mathbin{.}\lambda_{-1}(N)
  &=\lambda^p\bigl((x+y)\mathbin{.}\lambda_{-1}(N)\bigr)\\
  &=\sum_{i+j=p}\lambda^i(x\mathbin{.}\lambda_{-1}(N))\mathbin{.}\lambda^j(y\mathbin{.}\lambda_{-1}(N))\\
  &=\sum_{\substack{i+j=p\\ i,j\ne0}}\lambda^i(N,x)\mathbin{.}\lambda^j(N,y)\mathbin{.}(\lambda_{-1}(N))^2
    +(\lambda^p(N,x)+\lambda^p(N,y))\lambda_{-1}N
\end{align*}

Comme l'anneau considéré est intègre, on en déduit la relation suivante, valable par substitution dans tout @@M15@@-anneau, et pour tout élément @@M16@@ tel que @@M17@@ pour @@M18@@ :
@@M19@@

Considérons maintenant le @@M20@@-anneau engendré sur @@M21@@ par des indéterminées @@M22@@ soumises aux relations @@M23@@ pour @@M24@@, et @@M25@@ pour @@M26@@. On a alors @@M27@@ pour @@M28@@~; on peut donc définir les @@M29@@. On obtient :
@@M30@@
Par additivité de @@M31@@, on a : @@M32@@. On en déduit :
\begin{align*}
  \lambda^p(N+N',x)\mathbin{.}\lambda_{-1}(N)\mathbin{.}\lambda_{-1}(N')
  &=\lambda^p(x\mathbin{.}\lambda_{-1}(N)\mathbin{.}\lambda_{-1}(N'))\\
  &=\lambda^p(N',x\mathbin{.}\lambda_{-1}(N))\mathbin{.}\lambda_{-1}(N').
\end{align*}
